\documentclass{situationist-preprint}

\title{Universalizability and the Individuation of Maxims}
\msid{ETH-04-b}
\msfamily{ETH-04}
\msdiscipline{Normative and Applied Ethics}
\mstype{Companion manuscript}
\mskeywords{Kant; formula of universal law; maxims; individuation; false negatives; coordination maxims}

\begin{document}
\maketitle

\begin{abstract}
ETH-04-a argues that Kant's case against lying to the murderer at the door fails because it tests a maxim individuated too coarsely, and states a criterion on which a maxim includes exactly the circumstances the agent treats as reasons. We give independent support for the criterion. The formula of universal law is known to yield false prohibitions in cases of coordination and timing, such as playing tennis when the courts are empty. We show that these false prohibitions have the same source as the prohibition on the lie, coarse individuation, and that the criterion of ETH-04-a removes them while preserving the prohibition on the lying promise. A single diagnosis thus covers the murderer case and a family of standard objections to the formula.
\end{abstract}

\section{The criterion and its test}

ETH-04-a takes a maxim to have the form \enquote{in circumstances $C$, to do $A$, in order to bring about $E$} (Definition~2.1 of ETH-04-a), and individuates it by two principles: the circumstance-type $C$ includes every feature the agent treats as a reason for doing $A$ in pursuit of $E$, and no other (Principle~P1), and its features are general, specified without reference to particular persons, places or times (Principle~P2). It then shows that the coarse maxims of lying fail the contradiction in conception test (Lemma~3.1) while the householder's maxim passes both tests (Lemmas~3.2 and~3.3), so that the formula of universal law does not forbid the lie (Theorem~3.4), and the contrary argument tests a maxim that is too coarse (Corollary~3.5).

The problem of which description of an action is the relevant one is old \citep{ONeill1975}, and a criterion of individuation fitted to one case invites the suspicion that it was chosen to deliver that case. The suspicion is answered if the criterion, applied without adjustment, also corrects verdicts that are wrong for other reasons. We consider a class of such verdicts.

\section{Coordination maxims}

It is a familiar objection to the formula that it forbids innocent acts whose point depends on others not doing the same \citep{Herman1993,Parfit2011}. A player adopts the maxim \enquote{to play tennis on Sunday mornings, in order to find the courts empty}. If everyone played on Sunday mornings, the courts would not be empty then, and playing on Sunday mornings would not serve the end. On the practical reading of the contradiction in conception test, which ETH-04-a follows \citep{Korsgaard1985}, the maxim fails, and playing tennis on Sunday mornings is forbidden. That is absurd. Similar verdicts arise for buying a house when prices are low, travelling against the rush, and every maxim whose success depends on being in a minority.

\begin{proposition}\label{p:coord}
Individuated by Principles~P1 and~P2 of ETH-04-a, the tennis player's maxim passes the contradiction in conception test.
\end{proposition}

\begin{proof}
The player does not treat Sunday morning as a reason for playing; she treats the emptiness of the courts as the reason, and she plays on Sunday mornings because that is when she expects them to be empty. Her decision is sensitive to the emptiness, not to the day: if she learned that the courts were full on Sunday and empty on Tuesday, she would play on Tuesday. By Principle~P1 her maxim is \enquote{when the courts are uncrowded, to play tennis, in order to play on uncrowded courts}. In a world in which everyone acts on it, people play whenever the courts are uncrowded, and stop joining when they are crowded. Courts are uncrowded at some times in that world, and playing at those times achieves the end. The maxim is conceivable as a universal law and remains a means to its end.
\end{proof}

The coarse maxim fails for the same reason that the coarse maxims of lying fail in Lemma~3.1 of ETH-04-a: it substitutes for the feature the agent responds to a feature that merely correlates with it in her present circumstances. In the tennis case the substitution yields a false prohibition; in the murderer case it yields the false prohibition of the lie. In both, the corrected maxim includes the feature the agent's decision is sensitive to, and in both it passes.

\begin{proposition}\label{p:promise}
The criterion does not rescue the lying promise.
\end{proposition}

\begin{proof}
The agent who makes a lying promise to obtain money treats as reasons that she needs money and that a promise will produce it; she does not treat the identity of the lender, the day, or any further feature as a reason, since she would make the same promise to another lender on another day. By Principle~P1 her maxim is \enquote{when in need of money that a promise will produce, to promise repayment falsely, in order to obtain the money}. Universalized, promises made in need cease to be credible, and the false promise no longer produces the money; the maxim fails.
\end{proof}

The difference between Proposition~\ref{p:coord} and Proposition~\ref{p:promise} is principled. In the tennis case, the feature the agent responds to, emptiness, is one whose presence universal adoption adjusts to but does not abolish. In the promising case, the feature, the credibility of promises, is one that universal adoption destroys. The criterion of ETH-04-a does not decide which maxims pass; it decides which maxim is tested, and the test then does the work Kant assigned it.

\begin{objection}
An agent can always claim to respond to whatever feature would make her maxim pass. The criterion makes the test hostage to self-report.
\end{objection}

\begin{reply}
Principle~P1 concerns what the agent treats as a reason, and ETH-04-a (reply to Objection~1) makes this a matter of counterfactual sensitivity, not of avowal: a feature is in the maxim only if the decision would change were the feature absent. Whether the player would play on an empty Tuesday, or the false promiser would deceive a different creditor, is a fact about them, not something they can stipulate. The criterion is as objective as the agent's dispositions.
\end{reply}

\section{Conclusion}

The criterion of individuation of ETH-04-a corrects the false prohibitions generated by coordination maxims and preserves the prohibition of the lying promise. Its verdict on the murderer at the door is therefore not an ad hoc accommodation but an instance of a general correction: the formula of universal law, applied to maxims individuated by the agent's reasons, gives the right results in cases where coarse application fails.

\begin{thebibliography}{9}\small
\bibitem[Herman(1993)]{Herman1993} Herman, B. (1993). \emph{The Practice of Moral Judgment}. Harvard University Press.
\bibitem[Korsgaard(1985)]{Korsgaard1985} Korsgaard, C.~M. (1985). Kant's formula of universal law. \emph{Pacific Philosophical Quarterly}, 66(1--2), 24--47.
\bibitem[O'Neill(1975)]{ONeill1975} O'Neill, O. (1975). \emph{Acting on Principle: An Essay on Kantian Ethics}. Columbia University Press.
\bibitem[Parfit(2011)]{Parfit2011} Parfit, D. (2011). \emph{On What Matters}, Vol.~1. Oxford University Press.
\end{thebibliography}

\end{document}
